\section{Results}\label{sec:results}

% Every number, table and figure of this section comes from paper/generated/, which
% tools/paper_results.py writes from results/battery.csv, results/terms.csv, results/stats.json,
% results/ablation_triangulation.csv and paper/provenance.json (see paper/README.md). Predictions
% and the numbers of single diagnostic runs are taken from docs/history.md and
% docs/research/anomalies.md, are worded as such and are marked in comments.

\subsection{The runs, and the same law on three engines}\label{sec:runs}\label{sec:independence}

There are \result{battery/runs} runs: \result{battery/scenarios} scenarios,
\result{battery/methods} combinations of engine and contact method and up to
\result{battery/variants} variants per scenario. Of them, \result{battery/ok} finished with an
error value and \result{battery/failed} did not; the two laws have \result{status/laws/notok}
such run, and Section~S4.1 accounts for each. Errors are in percent of the reference, and the
tables show the base variant.

For \law{gonthier} and \law{distributed}, in the \result{ind/distributed/n} scenarios that
hold, the errors on the three engines differ by at most \result{independence/worst} (in
\result{ind/distributed/worst-scenario}), with a median relative difference of
\result{ind/distributed/median} (Section~S4.2). The three backends pass the same states to the
library and apply its wrenches at the same instant, on code bases that share nothing: this is a
check of the implementation, and the reason why one engine is enough to show a law.

\subsection{Class A: rigid-body references}\label{sec:classA}

Table~\ref{tab:errors-A} gives the errors on the \result{battery/n/A} class A scenarios. Where a
body slides or rolls on one plane (\scn{block}, \scn{cylinder}, \scn{sphere}, \scn{push},
\scn{vgroove}, \scn{twoblocks}, \scn{billiard}) both laws are within 1\,\%, except for the time at which rolling starts in \scn{billiard}
(\result{err/billiard/gonthier} and \result{err/billiard/distributed}), and they hold the
bodies that must be held (\scn{push}, \scn{wedge}, \scn{tower}). A patch model adds nothing
there: a rigid body on one plane is what point and complementarity contact are built for. The methods without a stick state creep long before
the friction limit: \law{gonthier-thesis} and \law{pcm} have \result{err/push/thesis} and
\result{err/push/pcm} in \scn{push}, and \result{err/wedge/thesis} and \result{err/wedge/pcm}
in \scn{wedge}. The larger errors of the two laws have four causes, each of which can be
predicted.

\begin{table}[t]
\caption{Class A: error [\%] of each method, base variant (abbreviations in
Section~\ref{sec:variants}; methods of the library on Chrono). The columns of the native models
are results of our configurations (Section~S3). \textit{nan}, \textit{timeout}: runs without an
error value}
\label{tab:errors-A}
\generatedtable{errors_class_A}
\end{table}

\emph{Rolling resistance of the damped foundation} (\scn{bowl}, \scn{trough}, \scn{turntable}).
In \scn{bowl} and \scn{trough} the period is right (for \law{distributed},
\result{term/bowl/distributed/period} and \result{term/trough/distributed/period}); the error is
the loss of amplitude, which the rigid theory keeps constant: \result{err/bowl/distributed} and
\result{err/trough/distributed}. In \scn{turntable}, whose rigid orbit is a closed circle, the
orbit rate is right (\result{term/turntable/distributed/orbit_rate}) and the orbit spirals
outwards: its radius grows by \result{err/turntable/distributed}. These errors are the rolling
resistance $\vect{M}_r = -K h\,\tens{J}\,\w_t$ of Eq.~\eqref{eq:gonthier-normal}. For the cap of
a sphere of radius $R$ pressed to a depth $\delta$, $\tens{J}$ has the tangential moment
$V R\delta/3$, so the torque is $-(h\,m g R\delta/3)\,\w_t$. In \scn{turntable} the no-slip
equations with that torque give a radius that grows as $\exp(\tfrac{25}{42}\,h g\,\delta\,t/R)$,
that is, by 77\,\% over the run for the static penetration of the scenario, against
\result{err/turntable/distributed} measured; in \scn{bowl} the same torque predicts a loss of
amplitude of 12\,\%, against \result{err/bowl/distributed} measured.
% 77 % and 12 % (0.351 rad left of 0.4): predictions derived in docs/history.md (entries of 2026-11-11), not battery results.
The errors are the model's answer at the $h$ and $K$ of the scenario, not numerical errors: any
damped Winkler foundation has this torque, every Winkler-type method of
Table~\ref{tab:errors-A} shows the loss, and with
$K\times10$ the errors of \law{distributed} fall from \result{err/turntable/distributed/dt0.2}
to \result{err/turntable/distributed/stiff} and from \result{err/bowl/distributed/dt0.2} to
\result{err/bowl/distributed/stiff}.

\emph{Compliance} (\scn{edge}). A rigid block leaves the table when its centre of mass passes
the edge; on a compliant foundation it pitches earlier, and every Winkler-type method does
(\result{err/edge/distributed} for the two laws, \result{err/edge/hydro} for Drake
hydroelastic). With $K\times10$ the error of the laws falls to
\result{err/edge/distributed/stiff}. These four scenarios measure $h$ and $K$ more than the
contact law, and summaries of class A are also given without them (``reduced'' class A).

\emph{The stability bound of the explicit coupling} (\scn{massratio}). The light block between
the heavy one and the ground is held by two contacts, so Eq.~\eqref{eq:stability} requires
$\dt < \sqrt{m/(K A_c)}$, which is 0.25\,ms for the scenario.
% K A_c = 1.6e7 N/m, m = 1 kg: design parameters of the scenario (docs/scenarios.md)
The two laws are lost at 2, 1 and 0.5\,ms and hold from 0.2\,ms, with
\result{err/massratio/distributed/dt0.2} of error at 0.2\,ms and
\result{err/massratio/distributed/dt0.1} at 0.1\,ms, a residual speed that the bounded dampers
remove slowly. A formulation that is implicit in the contact force has no such bound, and
complementarity contact holds the scenario at the base step. The tower of ten blocks is inside the
bound and both laws hold it at the base step (\result{err/tower/distributed}).

\emph{One normal per pair of convex pieces} (\scn{journal}). The error of \law{distributed}
(\result{err/journal/distributed}) is in the angle at which the shaft rides up the hole, and it
does not fall with the mesh (Section~\ref{sec:convergence}).

Two more results concern the laws. On a ground cut into 1, 2, 8 and 64 pieces
\law{distributed} gives the same error (\result{err/pieces64/distributed}); \law{gonthier} is
unaffected up to 8 pieces and turns the spinning box differently on 64
(\result{err/pieces64/gonthier}), and \law{pcm} depends on the cut
(\result{err/pieces64/pcm}). In \scn{oblique} the laws miss the rigid impulse theory by
\result{err/oblique/gonthier} and \result{err/oblique/distributed}.

\emph{Energy} (\scn{energy}). The error is the change of kinetic plus gravitational energy of a
frictionless, undamped ball between the first and the last period. The measure omits the energy
held by the undamped bounce on the contact spring, and sinking to the static penetration alone
moves 2.9\,\% of the energy of the swing, % 2.9 %: m g delta over the swing's energy, docs/research/anomalies.md §2
so it varies irregularly with the step and with the start for every law: over the five steps of
the series between \result{err/energy/gonthier/dt0.2} and \result{err/energy/gonthier/dt2} for
\law{gonthier}, between \result{err/energy/pcm} and \result{err/energy/pcm/dt0.2} for \law{pcm}
and between \result{err/energy/distributed/dt0.5} and \result{err/energy/distributed/dt0.1} for
\law{distributed}. The \result{err/energy/gonthier} of \law{gonthier} in
Table~\ref{tab:errors-A} is one value of that measure and not a steady gain. There is also a
true, smaller effect: Gonthier's normal force $K V\n$ is the gradient of a potential only where
the contact area is a function of the volume alone; with $\nabla V = A_c\n$ its curl is
$-K\,(V/A_c^2)\,\nabla A_c\times\nabla V$, which does not vanish on faceted or decomposed
surfaces. Its size was not separated from the noise of the measure (Section~S4.7).

\subsection{Class B: references that assume a pressure distribution}\label{sec:classB}

Table~\ref{tab:errors-B} gives the errors against the nominal reference and marks the position
of each run against the band of admissible distributions, which is the part of the reference
that does not depend on the Winkler model.

\begin{table}[t]
\caption{Class B: error [\%] against the nominal reference, base variant. Band: admissible
pressure distributions, in percent of the nominal stop time. $^\dagger$: the error exceeds the
wider side of the band, so the run is outside it; $^{*}$: it exceeds the narrower side only, so
its position is undetermined (the sign of the error was not recorded); unmarked errors are
inside. The columns of the native models are results of our configurations (Section~S3)}
\label{tab:errors-B}
\generatedtable{errors_B}
\end{table}

No elastic-foundation method is outside the band in any scenario: \law{distributed},
\law{gonthier}, \law{pcm} and Drake hydroelastic have no run marked as outside
(\law{gonthier-thesis} is the exception among the methods of the library). \law{distributed}
and Drake hydroelastic, the two methods that integrate the foundation over the contact surface,
are known to be inside the band in \result{band/distributed/inside} and
\result{band/hydro/inside} of the five scenarios and are undetermined in the others;
\law{gonthier} and \law{pcm} are inside in fewer. The two also meet the nominal value, within
\result{err/bearing/distributed} and \result{err/cone/hydro}. That agreement is a consistency
result: the nominal value is the model's own distribution, so two independent implementations
of the same pressure model, on different engines and with different discretizations, integrate
it correctly, and nothing follows about whether that pressure is the true one. How far it is
from the pressure of a continuum support is measured, for two of the scenarios, in
Section~\ref{sec:fem}.

\law{gonthier} misses the nominal value by the amount that its lumping predicts:
\result{err/disk/gonthier} in \scn{disk}, where the torque through the radius of gyration is
6\,\% above the exact integral (Section~\ref{sec:gonthier}), and
\result{err/contensou/gonthier} in \scn{contensou}. \law{pcm} has the error of its one-point
friction on the battery's meshes (\result{err/disk/pcm}; Section~\ref{sec:ablation}) and the
smallest error of all methods in \scn{bearing} (\result{err/bearing/pcm}), where its depth
follows the local surface.

A point-contact model does not represent the distribution. A face held by points of its
rim has the friction torque of a ring, $\mu N R$ against $\tfrac23\mu N R$ for the disk, unless
the user supplies a torsional friction coefficient for each geometry; the same holds for
complementarity contact on the same contact points. In the table, each of these models is
outside the band in \result{band/smc/outside} to \result{band/dpoint/outside} of the five
scenarios, and several of their runs scatter over nearby starts or end without a stop time
(Section~S3.4).

% An independent finite-element reference for two class B scenarios (disk, bearing).
% Input by sections/05-results.tex after the class B results. The solver, the cases and the
% numbers of Table~\ref{tab:fem} come from tools/fem (uv run tools/fem/run.py; results/fem/).
% The abstract, Section 7 and Section 8 repeat the 18 % of disk; change them with this file.
\subsection{A finite-element reference for two class B scenarios}\label{sec:fem}

The class B references assume a pressure distribution, and the band lists the alternatives. To
place a continuum in that band we solved, by finite elements, the static frictionless contact of
a linear-elastic support with a rigid body (small strain, quadratic quadrilaterals, unilateral
contact solved exactly by an active set; \code{tools/fem}). For \scn{disk}, an axisymmetric layer
of thickness $H$ bonded to a rigid base under a rigid flat punch of radius $R$; the torque of
sliding everywhere is $f\mu NR$ with $f=\int p\,r\,\dd A/(NR)$. The friction of a spin is
circumferential and, in linear elasticity, uncoupled from the normal problem, so the frictionless
pressure is exact for the spinning punch. For \scn{bearing}, a plane-strain sleeve of wall $t$
bonded to a rigid housing around a rigid journal without clearance, with $E$ such that
$E(1-\nu)/((1+\nu)(1-2\nu)t)$ is the battery's $K$; the torque with the friction circle is
$f\mu Nr/\sqrt{1+\mu^2}$ with $f=\int p\,\dd s/|\int p\,\n\,\dd s|$. There the friction couples
with the normal problem, and the frictionless pressure is its $\mu\to0$ limit. The solver
reproduces the flat punch on a deep layer ($f=\pi/4$ within 0.2\,\%), the Winkler limit of a thin
compressible layer ($f=2/3$) and Hertz line contact (half-width within one node, pressure within
1\,\% of its peak). Over three refinements the ratios of \scn{disk} converge with observed order
2.2 to 2.8 and those of \scn{bearing} change by less than $10^{-4}$; every digit of
Table~\ref{tab:fem} is converged; the pressures are plotted in Section~S4.4. The stop time follows from the scenario's own formula,
$t=I\omega_0/\tau\propto1/f$.

\emph{\scn{disk}.} A thin layer is a Winkler foundation: for $H/R\le0.2$ the pressure is nearly
uniform and $f$ is within 3\,\% of the nominal $2/3$ (below it for $\nu=0.45$, towards the
parabolic pressure of an incompressible thin layer). The shear between neighbouring fibres, which
a Winkler foundation omits, moves the load to the rim as the layer thickens, and the battery's
slab ($H/R=4$) is a half-space: $f$ is $\pi/4$ within 0.3\,\% for both $\nu$. A continuum support
puts \scn{disk} at the rigid-punch end of its band, with a stop time 15\,\% shorter than the
nominal one. No elastic-foundation method models that rim: \law{distributed} and Drake
hydroelastic, which reproduce the uniform pressure of their own model, stop 18\,\% late against
the FE reference, \law{pcm} 29\,\%; \law{gonthier} comes closest (11\,\%) because its lumped
torque is 6\,\% high, an error of the lumping and not a model of the rim. The 18\,\% is therefore a
limit of the Winkler foundation itself on a thick support, shared by its two independent
integrations, and not an error of their discretization; a finer mesh or a smaller step would not
reduce it.

\emph{\scn{bearing}.} For $\nu=0.3$ the continuum sleeve gives the Winkler cosine ($f$ within
0.2\,\% of $4/\pi$, contact over $\pm86^\circ$); a nearly incompressible wall ($\nu=0.45$)
spreads the load beyond the lower half ($\pm94^\circ$) and lies between the cosine and the
$\sqrt{\cos}$ distribution. The two ends of the band, line contact and uniform pressure, are
excluded. Against the FE stop time \law{pcm} and Drake hydroelastic are within 0.7\,\% and
\law{distributed} and \law{gonthier} within 2\,\% at $\nu=0.3$, and all four within 5\,\% at
$\nu=0.45$: the reference does not separate them, since its own dependence on $\nu$ is as large as
their differences, but it confirms that the nominal reference of \scn{bearing} is that of a
continuum and that the large errors of the point-contact models are theirs.

What the FE reference does not settle: it is static, small-strain and frictionless in the normal
problem, and it assumes sliding everywhere; the runs are dynamic and damped. The support's
$\nu$ is not given by the battery, and $E$ follows from $K$ through the thin-layer mapping, which
fixes the sink (the continuum's differs from the Winkler sink by $-5$ to $+22\,\%$) but not $f$.
The battery's ground is defined as a Winkler foundation, so agreement with the nominal reference
remains the right test of whether a method integrates its own model; the FE reference measures
how far that model is from a continuum. \scn{ring} and \scn{contensou}, flat faces on the same
slab, should move towards the rim as \scn{disk} does; they and \scn{cone} were not computed. The
signs in Table~\ref{tab:fem} come from the viewer's export of the committed battery, which
\code{tools/fem/run.py} checks against \code{results/battery.csv}.

% generated by tools/fem: typed from results/fem/fem_reference.csv (converged rows) and
% results/fem/stop_times.csv (engine Drake for the methods of the library; Chrono and MuJoCo
% differ by less than 0.1 %). Rerun tools/fem/run.py and retype if either file changes.
\begin{table}[!tbp]
\caption{Finite-element reference. Top: $f$ of a continuum support against the band (\scn{disk}:
nominal $2/3$, half-space punch $\pi/4\approx0.785$, uniform wear $1/2$; \scn{bearing}: nominal
$4/\pi\approx1.273$, line contact 1, $\sqrt{\cos}$ 1.371, uniform $\pi/2$). Bottom: deviation of
the base-variant stop time [\%] from the nominal reference and from the FE one (\scn{disk}: the
slab with $\nu=0.3$; $\nu=0.45$ moves it by at most 0.1 point)}
\label{tab:fem}
\centering\small
\begin{tabular}{@{}lcccc@{}}
\toprule
\scn{disk}, $f$ & $H/R=0.05$ & $0.2$ & $1$ & $4$ (battery) \\
\midrule
$\nu=0.3$  & 0.671 & 0.685 & 0.750 & 0.784 \\
$\nu=0.45$ & 0.654 & 0.659 & 0.741 & 0.783 \\
\midrule
\scn{bearing}, $f$ & small load & battery load & contact & $E$ [kPa] \\
\midrule
$\nu=0.3$  & 1.272 & 1.276 & $\pm86^\circ$ & 14.9 \\
$\nu=0.45$ & 1.309 & 1.314 & $\pm94^\circ$ & 5.27 \\
\bottomrule
\end{tabular}

\medskip
\begin{tabular}{@{}lrrrrr@{}}
\toprule
& \multicolumn{2}{c}{\scn{disk}} & \multicolumn{3}{c}{\scn{bearing}} \\
\cmidrule(lr){2-3}\cmidrule(l){4-6}
Method & nominal & FE & nominal & FE, $\nu=0.3$ & FE, $\nu=0.45$ \\
\midrule
\law{distributed}  & $+0.1$  & $+17.7$ & $+1.6$  & $+1.8$  & $+4.9$ \\
\law{gonthier}     & $-5.6$  & $+11.0$ & $-1.8$  & $-1.6$  & $+1.4$ \\
\law{pcm}          & $+10.0$ & $+29.3$ & $-0.2$  & $+0.0$  & $+3.1$ \\
Drake hydroelastic & $+0.3$  & $+18.0$ & $+0.5$  & $+0.6$  & $+3.7$ \\
\bottomrule
\end{tabular}
\end{table}


\subsection{Class C: consistency with the Winkler foundation}\label{sec:classC}

Class C asks whether a method is a Winkler foundation of stiffness $K$
(Table~\ref{tab:errors-C}); it does not apply to complementarity contact, which has no
compliance, or to a point contact whose stiffness is not areal. \law{distributed} is within
5\,\% in \result{pass/C/distributed/5} scenarios and \law{gonthier} in
\result{pass/C/gonthier/5}; the largest errors are the depth reached in the press fit
(\scn{hole}: \result{err/hole/gonthier} and \result{err/hole/distributed}). For the two laws
these scenarios are regression tests against the law's own closed-form predictions: the tilt
under an eccentric load, the tipping angle on a soft support and the depth of a press fit all
follow from the pressure over the patch. \law{pcm} and \law{gonthier-thesis} fail the scenarios
in which a body must be held by friction (\scn{tip}, \scn{hole}, \scn{stack}). Drake
hydroelastic, the same pressure model, is within 5\,\% in \result{pass/C/hydro/5}; its errors
in the impact scenarios fall with the step (\scn{drop}: \result{err/drop/hydro} at 1\,ms,
\result{err/drop/hydro/dt0.1} at 0.1\,ms), and three others do not and are unexplained
(Section~S3.5).

\begin{table}[t]
\caption{Class C: error [\%] of the methods that are a Winkler foundation, base variant}
\label{tab:errors-C}
\generatedtable{errors_class_C}
\end{table}

\subsection{The friction quadrature and the triangulation}\label{sec:ablation}

On the meshes of the battery the closed-form friction and the two quadratures of the
literature differ widely on the flat contacts that spin: in \scn{disk},
\result{err/disk/distributed} with the closed form against \result{err/disk/vertex} with the
vertex rule and \result{err/disk/centroid} with one point per triangle, at 128 segments per
turn as at 64 (\result{err/disk/vertex/mesh128} and \result{err/disk/centroid/mesh128}); and
with the ground cut into 64 pieces the error of \scn{pieces} goes from
\result{err/pieces1/vertex} to \result{err/pieces64/vertex} with the vertex rule and to
\result{err/pieces64/centroid} with the one-point rule, while the closed form stays at
\result{err/pieces64/distributed} (Section~S4.5). Where nothing spins about the normal of a
face the three rules agree.

The size of those differences is a property of the battery's meshes. The flat cap of a disk is
triangulated as a fan from a vertex of its rim: long triangles that span the face, none with a
vertex at the centre of rotation, and more segments per turn only add more of them. To separate
the rule from the mesh, the three rules were evaluated, in a re-implementation apart from the
library, on a disk and a ring under a uniform pressure spinning in place, and under a tilted
pressure sliding while spinning about a point at half the radius, for several triangulations of
the same face (Table~\ref{tab:triangulation}, Fig.~\ref{fig:triangulation}). The closed form
gives the same wrench on every triangulation, to round-off (largest difference
\result{tri/closed-spread}). On the fan from the rim the one-point rule loses
\result{tri/disk/spin/rim-fan/64/centroid/moment} of the spin torque and the vertex rule
\result{tri/disk/spin/rim-fan/64/vertex/moment}, at any number of segments: these are the
errors behind the battery's stop times. On a fan from the centre the one-point rule has
\result{tri/disk/spin/centre-fan/64/centroid/moment}, while the vertex rule, which interpolates
unit vectors from a vertex that does not slip, still loses
\result{tri/disk/spin/centre-fan/64/vertex/moment}; when the face also slides the centre fan
is no longer a lucky mesh for the one-point rule
(\result{tri/disk/slide+spin/centre-fan/64/centroid/force} of the load in the force). On
well-shaped triangles both rules converge at second order: with
\result{tri/disk/spin/hex/4/centroid/triangles} nearly equilateral triangles the errors are
\result{tri/disk/spin/hex/4/centroid/moment} and \result{tri/disk/spin/hex/4/vertex/moment},
and with \result{tri/disk/spin/hex/16/centroid/triangles},
\result{tri/disk/spin/hex/16/centroid/moment} and \result{tri/disk/spin/hex/16/vertex/moment}.

So the closed form is not more accurate than a well-used quadrature: on a face meshed with a
few hundred well-shaped triangles a one-point rule is within a few tenths of a percent, the
level of the other errors of the battery, and Drake hydroelastic, which uses one point per
polygon of a refined contact surface, has \result{err/disk/hydro} in \scn{disk}. What the
closed form gives is independence: the friction wrench of a flat face does not depend on how
the face is triangulated or on how the bodies are cut into pieces, so a flat face needs no mesh
of its own. A quadrature needs the face remeshed, and its error then depends on the quality of
that mesh.

\begin{table}[t]
\caption{The friction quadratures on one flat face, apart from the battery: error [\%] of the
friction torque of a face spinning in place under uniform pressure, and of the torque and the
force (the latter over $\mu$ times the load) of a face sliding and spinning under a tilted
pressure, against the closed form, which is the same on every triangulation. Negative: too
little torque}
\label{tab:triangulation}
\generatedtable{triangulation}
\end{table}

\begin{figure}[t]
\centering
\generatedfigure[width=0.9\linewidth]{ablation_triangulation}
\caption{Error of the spin torque of a disk under uniform pressure with the one-point rule and
with the vertex rule, against the number of triangles of the face, for five families of
triangulations. The closed form has no error on any of them}
\label{fig:triangulation}
\end{figure}

\subsection{Convergence with the step and the mesh}\label{sec:convergence}

\emph{Step.} The step series of the two laws are mostly flat (a slope of the logarithm of the
error within $\pm0.2$, or every error below $2\times10^{-3}$): \result{convn/dt/distributed/flat}
of \result{convn/dt/distributed/n} for \law{distributed} and \result{convn/dt/gonthier/flat} for
\law{gonthier} (Section~S4.6). Their error at 1\,ms is then mostly not an error of the time
integration but the error of the model against the reference, or of the mesh. Where it is a
step error it falls roughly in proportion to the step (\scn{drop}, \scn{tower}), and
\scn{massratio} shows the stability bound (Section~\ref{sec:classA}).

\emph{Mesh.} The errors of \law{distributed} fall with the number of segments in
\result{convn/mesh/distributed/converging} of the \result{convn/mesh/distributed/n} scenarios
with curved surfaces. Two series that do not converge matter. \scn{journal} is flat
(\result{err/journal/distributed/mesh32}, \result{err/journal/distributed} and
\result{err/journal/distributed/mesh128} at 32, 64 and 128 segments), while \law{pcm} converges
on it (\result{err/journal/pcm/mesh128} at 128): the geometry of \law{distributed} is exact for
the polyhedra, but with one normal per pair of pieces its answer for a curved conforming
contact does not converge to that of the curved bodies. \scn{cone} is not monotone
(\result{err/cone/distributed/mesh32}, \result{err/cone/distributed},
\result{err/cone/distributed/mesh128}), at errors below 1\,\%. In \scn{bearing} the error of
\law{distributed} falls from \result{err/bearing/distributed/mesh32} to
\result{err/bearing/distributed/mesh128}, and that of Drake hydroelastic from
\result{err/bearing/hydro/mesh32} to \result{err/bearing/hydro/mesh128}: a contact surface
without a representative normal converges faster and further on curved conforming contact. The
series have three points, so no order is quoted.

\subsection{A stiffer foundation: friction-induced flutter}\label{sec:stiffness}

The stiffness series has two points, $K$ and $10K$, both at 0.2\,ms: enough to see a direction
and not a limit (Section~S4.7). On class A the error of \law{distributed} falls by more than a
tenth in \result{stiffn/distributed/better} of \result{stiffn/distributed/n} scenarios, among
them those whose error was attributed above to the compliance or to the damping of the
foundation, and
rises in \result{stiffn/distributed/worse} (\scn{cylinder}, and \scn{massratio}, which is
further beyond the stability bound). On class B the stiffer foundation makes the errors of the
two laws larger in \result{stiffn/distributed/worse/B} of the five scenarios: \scn{disk} from
\result{err/disk/distributed/dt0.2} to \result{err/disk/distributed/stiff}, \scn{cone} from
\result{err/cone/distributed/dt0.2} to \result{err/cone/distributed/stiff}, \scn{bearing} from
\result{err/bearing/distributed/dt0.2} to \result{err/bearing/distributed/stiff}.

On the flat \scn{disk} the cause is not the explicit coupling: in diagnostic runs the error at
$10K$ was the same at steps from 1 to 0.05\,ms.
% Diagnostic single runs on Chrono, docs/research/anomalies.md §1: 18.1 % at 0.2 ms, 17.7 % at 0.05 ms; tilt 1.8e-5 rad at K, 1.2e-3 rad at 10 K; growth 6.6 and 32 1/s measured, 6.3 and 31 predicted.
It is a self-excited wobble of the model. A tilt $\vect{\theta}$ of the spinning disk shifts
the pressure, so the friction of the heavier side no longer cancels that of the lighter one: a
net friction force of magnitude $\mu K\,(\pi R^3/3)\,\theta$ appears, and it acts at the
contact, half the thickness $H$ of the disk below the centre of mass. Its moment,
$N_c = \mu K\,(\pi R^3/3)\,(H/2)$ per unit of tilt, turns the tilt instead of restoring it (a
circulatory term), and the whirl grows when
\begin{equation}
  \mu K\,\frac{\pi R^3}{3}\,\frac{H}{2} \;>\; c_\theta\,\omega_\theta ,
  \qquad c_\theta = \frac{h\,m g R^2}{4}, \qquad \omega_\theta = \sqrt{\frac{K I_A}{I_t}} ,
  \label{eq:flutter}
\end{equation}
with $c_\theta$ the rocking damping of the foundation under the weight $m g$,
$I_A = \pi R^4/4$ the second moment of the face and $I_t$ the transverse moment of inertia of
the disk, at the rate $(N_c/\omega_\theta - c_\theta)/(2 I_t)$. For the disk of the scenario
the threshold is $K = 1.9\times10^6$\,N/m$^3$, about twice the scenario's stiffness: at $K$ the
wobble decays, at $3K$ the condition predicts a growth of 6\,s$^{-1}$ and the diagnostic run
showed 6.6, and at $10K$ it predicts 31\,s$^{-1}$ against 32 measured. The tilt saturates at
$1.2\times10^{-3}$\,rad; the disk then rides on its rim, where the friction torque tends to
$\mu N R$, and stops early. The left side of Eq.~\eqref{eq:flutter} grows with $K$ and the
right side with $\sqrt{K}$: a foundation made stiffer at the same $h$ is less damped in
rocking. In general the condition is $\mu K S_\rho\,z_c > c_\theta\,\omega_\theta$, with
$S_\rho$ the first radial moment of the patch and $z_c$ the height of the centre of mass.
Eq.~\eqref{eq:flutter} is therefore a limit of validity of the damped Winkler law with sliding
friction under a body whose centre of mass is above the contact, and not a defect of an
implementation: the linear model reproduces the growth rates of both laws. On the curved seats
(\scn{cone}, \scn{bearing}) every Winkler-type method degrades at $10K$, Drake hydroelastic
included (\result{err/cone/hydro/dt0.2} to \result{err/cone/hydro/stiff}); there the same
self-excited motion and the sag of the facets, which is no longer small against a penetration
ten times smaller, act together and were not separated.

\subsection{Cost}\label{sec:cost}

% The cost numbers are generated from results/timing.csv (scripts/bench.sh --timing on an idle
% machine); without that file the generator falls back to the accuracy pass and labels it.

Costs are wall-clock time over simulated time from \result{cost/source}: every scenario in the
base variant, one run at a time, the median of \result{battery/timing-repeats} repeats, on an
\result{battery/timing-machine} processor (Section~S4.8). \law{gonthier} runs at a median of
\result{cost/gonthier/median} of real time and never above \result{cost/gonthier/max}.
\law{distributed} has a median of \result{cost/distributed/median} and is slower than real time
in \result{cost/distributed/slow} scenarios, the two with spheres hitting a plane, up to
\result{cost/distributed/max}; its cost grows with the number of faces in contact. \law{pcm}
has a median of \result{cost/pcm/median}. The point-contact and complementarity models, as
configured, have medians between \result{cost/mjdefault/median} and
\result{cost/dpoint/median}. Drake hydroelastic computes a contact surface between tetrahedral
pressure fields wherever the ground is not a plane, and solves one constraint per polygon of
that surface; as configured here and at the step of the battery, it has a median of
\result{cost/hydro/median}, is slower than real time in \result{cost/hydro/slow} of the
\result{cost/hydro/n} scenarios that it finished, up to \result{cost/hydro/max} times, and did
not finish \scn{tower} within the limit.

\subsection{Paired comparisons}\label{sec:ratios}

Across scenarios, two methods are compared by the geometric mean of the ratio of their errors
on the same scenarios, with errors capped at \result{rules/cap}, runs without a value counted
at the cap, a bootstrap interval and Holm's correction over a family of \result{primary/n}
primary comparisons: \law{distributed} against \law{gonthier}, and each of the two against one
native model per engine, on class A and on classes A and B (Table~\ref{tab:primary};
Section~S5 has the rules, the sensitivity analyses and the exploratory pairs). Only
\result{primary/significant} of the \result{primary/n} survive the correction, both on classes
A and B together: \law{distributed} has about half the error of \law{gonthier}
(\result{robust/distributed-vs-gonthier/A+B/full}, adjusted $p = $
\result{robust/distributed-vs-gonthier/A+B/holm}) and about a quarter of that of Chrono SMC
(\result{robust/distributed-vs-smc/A+B/full}), the latter narrowly (adjusted $p = $
\result{robust/distributed-vs-smc/A+B/holm}). On class A alone no difference is significant:
the ratios of the laws to the native models are below 1 with intervals that reach above it.
Within 5\,\% of the reference, \law{distributed} solves \result{pass/A/distributed/5} class A
scenarios and \result{pass/B/distributed/5} of class B, and \law{gonthier}
\result{pass/A/gonthier/5} and \result{pass/B/gonthier/5}. The scenarios are few, hand-picked
and correlated, so these figures describe this battery and are not an estimate for other
problems.

\begin{table}[t]
\caption{The primary comparisons: geometric mean of the ratio of the errors of two methods on
the same scenarios (base variant; a law against a native model on that model's engine), with
its 95\,\% bootstrap interval and the Holm-adjusted $p$ over the \result{primary/n}
comparisons. Bold: adjusted $p$ below 0.05}
\label{tab:primary}
\generatedtable{primary_compact}
\end{table}
